\section{Results}
\label{sec:results}
The results are mixed.
We notice intervals that have correct empirical coverage, but we also see intervals with incorrect empirical coverage.
Among those with correct empirical coverage, we see intervals that are statistically indistinguishable from exactly correct, but we also see intervals that are conservatively correct.
By ``statistically indistinguishable from exactly correct'' we mean the situation where the bootstrap/Wilson 95\% confidence interval on the empirical coverage estimate contains the nominal coverage.
We classify an interval as having incorrect empirical coverage when the 95\% confidence interval on the empirical coverage estimate lies completely below the nominal coverage.
Similarly, we classify an interval as having conservatively correct coverage when the 95\% confidence interval on the empirical coverage estimate lies completely above the nominal coverage.\footnote{Technically, for every claim of incorrect, exactly correct, or conservatively correct coverage we should add the words ``statistically indistinguishable from''. For brevity we will omit these words from now on.}

We show two tables here: \cref{tab:coverage} on the empirical coverage of the intervals on $\theta_r$ averaged over $r$, and \cref{tab:coverage_Eb1} on the empirical coverage of the intervals on $\mathbb{E}[\beta_1]$.
The empirical coverage of intervals on $\mathbb{E}[\beta_2]$ can be found in \cref{tab:coverage_Eb2} in \cref{sec:app:tab}.
It is excluded here because it is very similar to \cref{tab:coverage_Eb1}.
We include clarifying figures in \cref{sec:app:fig}.
This includes figures of a single estimated PMF/CDF for both the OLS and FKRB estimators, for all DGPs, and it includes figures that specify the empirical coverage of individual $\theta_r$ over the whole grid.
These coverage plots allow us to identify any spots in the grid where the empirical coverage deviates from other points.

For the FKRB intervals for $\theta_r$ we see that the coverage is exactly correct on average for practically all off-boundary DGPs that have $R=25$.
It does not seem to depend much on the nature of the off-boundary DGP; near-boundary and far-from-boundary random coefficient distributions show similar behavior.
For the off-boundary DGPs with $R=225$ we notice something interesting: the FKRB intervals have incorrect coverage on average.
We suspect this stems from the OLS nature of the FKRB interval.
We do not know fully why, but the OLS estimator consistently has trouble estimating $\theta_r$ around the point $(-3,1)$.
See for example \cref{fig:olsnature}.
This is consistent across all distributions; see \cref{sec:app:fig} to see that this behavior is seen for all DGPs with $R=225$.
It is therefore unlikely that any distributional properties play a role in this.
It is more likely that the design matrices in the OLS estimator become ill-conditioned in a densely correlated grid. 
As a result of this behavior, the empirical coverage is extremely low around the point $(-3,1)$.
This skews the average result displayed in \cref{tab:coverage}.
As can also be seen in \cref{fig:olsnature}, the FKRB interval seems to behave fine in the other part of the grid.
The behavior is similar for both $\alpha=0.05$ and $\alpha=0.5$.

\begin{figure}[htbp]
    \centering
    \begin{minipage}{0.48\textwidth}
        \centering
        \includegraphics[width=\linewidth]{../output/figures/estimated/pmf_ols_beta_bimodal_N10000_R225.png}
    \end{minipage}\hfill
    \begin{minipage}{0.48\textwidth}
        \centering
        \includegraphics[width=\linewidth]{../output/figures/coverage/coverage_beta_bimodal_support_probs_alpha0p05_N10000_R225_fkrb.png}
    \end{minipage}
    \caption{Estimated OLS PMF of the bimodal DGP ($N=10{,}000$, $R=225$) on the left, with the corresponding grid plot of the empirical coverage of the FKRB interval at $\alpha=0.05$ on the right.}
    \label{fig:olsnature}
\end{figure}

The sieve Wald intervals for $\theta_r$ are highly conservative.
For $R=225$, not a single DGP generates an interval with an empirical coverage of less than 100\%.
For $R=25$ it is slightly better, but still highly conservative.
We do notice that boundary issues play a predominant role in the conservativeness.
For example, for the far-from-boundary strictly uniform DGP with $N=10{,}000$, $R=25$, $\alpha=0.5$, we see an empirical coverage of $0.656$.
This contrasts sharply with the exactly-on-boundary single type DGP with $N=10{,}000$, $R=25$, $\alpha=0.5$, which has an empirical coverage of $0.969$.
In \cref{fig:waldcovergage} we see that the interval is more conservative on grid points which are close to the boundary.
We also make an interesting observation on the single type DGP, where each grid point lies exactly on the boundary: the grid point that has a true probability mass of $1$ has a corresponding confidence interval that is not conservative.
We would expect that grid points close to $1$ have the same problems as grid points close to $0$, but this observation suggests otherwise.

\begin{figure}[htbp]
    \centering
    \begin{minipage}{0.48\textwidth}
        \centering
        \includegraphics[width=\linewidth]{../output/figures/coverage/coverage_beta_single_type_support_probs_alpha0p50_N10000_R25_cp_wald.png}
    \end{minipage}\hfill
    \begin{minipage}{0.48\textwidth}
        \centering
        \includegraphics[width=\linewidth]{../output/figures/coverage/coverage_beta_bimodal_support_probs_alpha0p50_N10000_R25_cp_wald.png}
    \end{minipage}
    \caption{Empirical coverage plots for the sieve Wald intervals of the single type and bimodal DGPs ($N=10{,}000$, $R=25$, $\alpha=0.5$).}
    \label{fig:waldcovergage}
\end{figure}

In general, the QLR intervals for $\theta_r$ behave similarly to the sieve Wald intervals, in the sense that they are quite conservative, although considerably less so.
They also exhibit the phenomenon where points with a probability mass equal or close to $1$ have a corresponding interval that is incorrect.
See for example \cref{fig:QLRcoverage}.
Interestingly, DGPs with more boundary issues are less conservative than DGPs with fewer boundary issues.
This possibly arises because, for these DGPs, grid points with positive probability mass have incorrect empirical coverage, which skews the average down.

\begin{figure}[htbp]
    \centering
    \begin{minipage}{0.48\textwidth}
        \centering
        \includegraphics[width=\linewidth]{../output/figures/coverage/coverage_beta_single_type_support_probs_alpha0p50_N10000_R225_qlr.png}
    \end{minipage}\hfill
    \begin{minipage}{0.48\textwidth}
        \centering
        \includegraphics[width=\linewidth]{../output/figures/coverage/coverage_beta_four_types_support_probs_alpha0p50_N10000_R225_qlr.png}
    \end{minipage}
    \caption{Empirical coverage plots for the QLR intervals of the single type and four types DGPs ($N=10{,}000$, $R=225$, $\alpha=0.5$).}
    \label{fig:QLRcoverage}
\end{figure}

For the intervals on $\mathbb{E}[\beta_1]$ and $\mathbb{E}[\beta_2]$ we see results similar to those for the $\theta_r$ intervals.
The uncertainty in the empirical coverage estimates is somewhat larger.
The FKRB intervals seem somewhat more robust across grid sizes, but overall they remain statistically indistinguishable from exactly correct.
Some DGPs have slightly lower empirical coverage than nominal, others slightly higher.
DGPs with more boundary issues are generally the ones with higher empirical coverage.

The sieve Wald intervals on the means are again highly conservative.
In fact, they are even more conservative than the ones on the individual $\theta_r$, with every interval at the $\alpha=0.05$ significance level having an empirical coverage of $1$.

The QLR intervals on the means are generally still conservative, but less so than the QLR intervals for the individual $\theta_r$.
We again see that the QLR intervals are less conservative for DGPs with a high degree of boundary problems.

\begin{landscape}
\begin{table}[htbp]
    \centering
    \caption{Empirical coverage of the original FKRB \citep{FoxKimRyanBajari2011}, and sieve Wald and sieve QLR \citep{ChenPouzo2015} confidence intervals for the support-point probabilities $\theta_r$, by data generating process, sample size $N$, grid size $R$ and significance level $\alpha$.}
    \label{tab:coverage}
    \footnotesize
    \setlength{\tabcolsep}{1.5pt}
    \begin{tabularx}{\linewidth}{ccc *{21}{>{\centering\arraybackslash}X}}
        \toprule
        & & & \multicolumn{3}{c}{Bimodal} & \multicolumn{3}{c}{Tight normal} & \multicolumn{3}{c}{Wide normal} & \multicolumn{3}{c}{Strictly uniform} & \multicolumn{3}{c}{Almost single type} & \multicolumn{3}{c}{Single type} & \multicolumn{3}{c}{Four types} \\
        \cmidrule(lr){4-6} \cmidrule(lr){7-9} \cmidrule(lr){10-12} \cmidrule(lr){13-15} \cmidrule(lr){16-18} \cmidrule(lr){19-21} \cmidrule(lr){22-24}
        $\alpha$ & $N$ & $R$ & FKRB & Wald & QLR & FKRB & Wald & QLR & FKRB & Wald & QLR & FKRB & Wald & QLR & FKRB & Wald & QLR & FKRB & Wald & QLR & FKRB & Wald & QLR \\
        \midrule
        0.05 & 1{,}000 & 25 & 0.951 & 0.997$^{*}$ & 0.979$^{*}$ & 0.944 & 1.000$^{*}$ & 0.990$^{*}$ & 0.942 & 1.000$^{*}$ & 0.992$^{*}$ & 0.944 & 0.999$^{*}$ & 0.992$^{*}$ & 0.946 & 1.000$^{*}$ & 0.996$^{*}$ & 0.976$^{*}$ & 1.000$^{*}$ & 0.984$^{*}$ & 0.966$^{*}$ & 0.996$^{*}$ & 0.957$^{*}$ \\
        & & & {\tiny [.944,~.958]} & {\tiny [.996,~.999]} & {\tiny [.975,~.983]} & {\tiny [.936,~.952]} & {\tiny [.999,~1.000]} & {\tiny [.988,~.993]} & {\tiny [.933,~.951]} & {\tiny [.999,~1.000]} & {\tiny [.989,~.994]} & {\tiny [.936,~.951]} & {\tiny [.999,~1.000]} & {\tiny [.989,~.994]} & {\tiny [.939,~.954]} & {\tiny [1.000,~1.000]} & {\tiny [.994,~.997]} & {\tiny [.972,~.980]} & {\tiny [1.000,~1.000]} & {\tiny [.981,~.987]} & {\tiny [.960,~.972]} & {\tiny [.995,~.998]} & {\tiny [.951,~.963]} \\[0.4ex]
        0.05 & 1{,}000 & 225 & 0.909$^{\dagger}$ & 1.000$^{*}$ & 1.000$^{*}$ & 0.900$^{\dagger}$ & 1.000$^{*}$ & 1.000$^{*}$ & 0.909$^{\dagger}$ & 1.000$^{*}$ & 1.000$^{*}$ & 0.908$^{\dagger}$ & 1.000$^{*}$ & 1.000$^{*}$ & 0.908$^{\dagger}$ & 1.000$^{*}$ & 0.999$^{*}$ & 0.953 & 1.000$^{*}$ & 0.997$^{*}$ & 0.955$^{*}$ & 1.000$^{*}$ & 0.992$^{*}$ \\
        & & & {\tiny [.903,~.915]} & {\tiny [1.000,~1.000]} & {\tiny [1.000,~1.000]} & {\tiny [.894,~.906]} & {\tiny [1.000,~1.000]} & {\tiny [1.000,~1.000]} & {\tiny [.903,~.914]} & {\tiny [1.000,~1.000]} & {\tiny [1.000,~1.000]} & {\tiny [.902,~.914]} & {\tiny [1.000,~1.000]} & {\tiny [1.000,~1.000]} & {\tiny [.902,~.913]} & {\tiny [1.000,~1.000]} & {\tiny [.999,~.999]} & {\tiny [.950,~.956]} & {\tiny [1.000,~1.000]} & {\tiny [.997,~.997]} & {\tiny [.952,~.958]} & {\tiny [1.000,~1.000]} & {\tiny [.992,~.993]} \\[0.4ex]
        0.05 & 10{,}000 & 25 & 0.952 & 0.992$^{*}$ & 0.969$^{*}$ & 0.947 & 0.998$^{*}$ & 0.977$^{*}$ & 0.949 & 0.992$^{*}$ & 0.970$^{*}$ & 0.950 & 0.991$^{*}$ & 0.970$^{*}$ & 0.955 & 1.000$^{*}$ & 0.955 & 0.973$^{*}$ & 0.999$^{*}$ & 0.983$^{*}$ & 0.969$^{*}$ & 0.995$^{*}$ & 0.968$^{*}$ \\
        & & & {\tiny [.945,~.958]} & {\tiny [.988,~.994]} & {\tiny [.963,~.974]} & {\tiny [.939,~.955]} & {\tiny [.996,~.999]} & {\tiny [.972,~.981]} & {\tiny [.942,~.957]} & {\tiny [.990,~.995]} & {\tiny [.964,~.975]} & {\tiny [.943,~.957]} & {\tiny [.988,~.993]} & {\tiny [.965,~.975]} & {\tiny [.947,~.962]} & {\tiny [1.000,~1.000]} & {\tiny [.944,~.965]} & {\tiny [.969,~.978]} & {\tiny [.998,~1.000]} & {\tiny [.980,~.986]} & {\tiny [.963,~.975]} & {\tiny [.993,~.997]} & {\tiny [.963,~.973]} \\[0.4ex]
        0.05 & 10{,}000 & 225 & 0.917$^{\dagger}$ & 1.000$^{*}$ & 1.000$^{*}$ & 0.922$^{\dagger}$ & 1.000$^{*}$ & 1.000$^{*}$ & 0.925$^{\dagger}$ & 1.000$^{*}$ & 1.000$^{*}$ & 0.939$^{\dagger}$ & 1.000$^{*}$ & 1.000$^{*}$ & 0.924$^{\dagger}$ & 1.000$^{*}$ & 1.000$^{*}$ & 0.960$^{*}$ & 1.000$^{*}$ & 0.997$^{*}$ & 0.957$^{*}$ & 1.000$^{*}$ & 0.993$^{*}$ \\
        & & & {\tiny [.911,~.922]} & {\tiny [1.000,~1.000]} & {\tiny [1.000,~1.000]} & {\tiny [.917,~.927]} & {\tiny [1.000,~1.000]} & {\tiny [1.000,~1.000]} & {\tiny [.919,~.930]} & {\tiny [1.000,~1.000]} & {\tiny [1.000,~1.000]} & {\tiny [.934,~.944]} & {\tiny [1.000,~1.000]} & {\tiny [1.000,~1.000]} & {\tiny [.919,~.930]} & {\tiny [1.000,~1.000]} & {\tiny [.999,~1.000]} & {\tiny [.957,~.963]} & {\tiny [1.000,~1.000]} & {\tiny [.997,~.998]} & {\tiny [.955,~.960]} & {\tiny [1.000,~1.000]} & {\tiny [.993,~.994]} \\[0.4ex]
        \cmidrule(l){1-24}
        0.5 & 1{,}000 & 25 & 0.479$^{\dagger}$ & 0.899$^{*}$ & 0.714$^{*}$ & 0.485 & 0.948$^{*}$ & 0.805$^{*}$ & 0.494 & 0.909$^{*}$ & 0.668$^{*}$ & 0.495 & 0.916$^{*}$ & 0.601$^{*}$ & 0.504 & 0.985$^{*}$ & 0.872$^{*}$ & 0.747$^{*}$ & 0.983$^{*}$ & 0.934$^{*}$ & 0.704$^{*}$ & 0.887$^{*}$ & 0.765$^{*}$ \\
        & & & {\tiny [.464,~.495]} & {\tiny [.890,~.908]} & {\tiny [.703,~.726]} & {\tiny [.468,~.502]} & {\tiny [.941,~.954]} & {\tiny [.794,~.815]} & {\tiny [.478,~.511]} & {\tiny [.901,~.917]} & {\tiny [.654,~.682]} & {\tiny [.478,~.512]} & {\tiny [.910,~.922]} & {\tiny [.585,~.617]} & {\tiny [.489,~.518]} & {\tiny [.982,~.989]} & {\tiny [.861,~.882]} & {\tiny [.738,~.756]} & {\tiny [.980,~.986]} & {\tiny [.929,~.938]} & {\tiny [.691,~.716]} & {\tiny [.878,~.896]} & {\tiny [.755,~.776]} \\[0.4ex]
        0.5 & 1{,}000 & 225 & 0.455$^{\dagger}$ & 1.000$^{*}$ & 0.993$^{*}$ & 0.454$^{\dagger}$ & 1.000$^{*}$ & 0.988$^{*}$ & 0.454$^{\dagger}$ & 1.000$^{*}$ & 0.996$^{*}$ & 0.452$^{\dagger}$ & 1.000$^{*}$ & 0.996$^{*}$ & 0.462$^{\dagger}$ & 1.000$^{*}$ & 0.995$^{*}$ & 0.728$^{*}$ & 1.000$^{*}$ & 0.994$^{*}$ & 0.726$^{*}$ & 1.000$^{*}$ & 0.980$^{*}$ \\
        & & & {\tiny [.446,~.464]} & {\tiny [1.000,~1.000]} & {\tiny [.992,~.994]} & {\tiny [.446,~.463]} & {\tiny [1.000,~1.000]} & {\tiny [.987,~.989]} & {\tiny [.445,~.462]} & {\tiny [1.000,~1.000]} & {\tiny [.995,~.996]} & {\tiny [.442,~.460]} & {\tiny [1.000,~1.000]} & {\tiny [.995,~.997]} & {\tiny [.453,~.471]} & {\tiny [1.000,~1.000]} & {\tiny [.994,~.995]} & {\tiny [.724,~.732]} & {\tiny [1.000,~1.000]} & {\tiny [.994,~.995]} & {\tiny [.721,~.731]} & {\tiny [1.000,~1.000]} & {\tiny [.979,~.980]} \\[0.4ex]
        0.5 & 10{,}000 & 25 & 0.502 & 0.785$^{*}$ & 0.709$^{*}$ & 0.489 & 0.887$^{*}$ & 0.757$^{*}$ & 0.493 & 0.738$^{*}$ & 0.628$^{*}$ & 0.502 & 0.656$^{*}$ & 0.532$^{*}$ & 0.500 & 0.960$^{*}$ & 0.750$^{*}$ & 0.748$^{*}$ & 0.969$^{*}$ & 0.950$^{*}$ & 0.709$^{*}$ & 0.857$^{*}$ & 0.792$^{*}$ \\
        & & & {\tiny [.485,~.519]} & {\tiny [.772,~.798]} & {\tiny [.694,~.724]} & {\tiny [.473,~.505]} & {\tiny [.879,~.895]} & {\tiny [.742,~.773]} & {\tiny [.476,~.510]} & {\tiny [.725,~.750]} & {\tiny [.609,~.648]} & {\tiny [.486,~.518]} & {\tiny [.644,~.669]} & {\tiny [.517,~.548]} & {\tiny [.483,~.516]} & {\tiny [.955,~.964]} & {\tiny [.734,~.767]} & {\tiny [.739,~.758]} & {\tiny [.966,~.972]} & {\tiny [.944,~.955]} & {\tiny [.699,~.720]} & {\tiny [.847,~.866]} & {\tiny [.780,~.804]} \\[0.4ex]
        0.5 & 10{,}000 & 225 & 0.476$^{\dagger}$ & 1.000$^{*}$ & 0.980$^{*}$ & 0.479$^{\dagger}$ & 1.000$^{*}$ & 0.980$^{*}$ & 0.477$^{\dagger}$ & 1.000$^{*}$ & 0.980$^{*}$ & 0.487$^{\dagger}$ & 1.000$^{*}$ & 0.979$^{*}$ & 0.477$^{\dagger}$ & 1.000$^{*}$ & 0.995$^{*}$ & 0.736$^{*}$ & 1.000$^{*}$ & 0.995$^{*}$ & 0.730$^{*}$ & 1.000$^{*}$ & 0.979$^{*}$ \\
        & & & {\tiny [.466,~.486]} & {\tiny [1.000,~1.000]} & {\tiny [.977,~.983]} & {\tiny [.471,~.487]} & {\tiny [1.000,~1.000]} & {\tiny [.978,~.982]} & {\tiny [.468,~.486]} & {\tiny [1.000,~1.000]} & {\tiny [.975,~.984]} & {\tiny [.478,~.496]} & {\tiny [1.000,~1.000]} & {\tiny [.972,~.984]} & {\tiny [.469,~.486]} & {\tiny [1.000,~1.000]} & {\tiny [.995,~.996]} & {\tiny [.731,~.741]} & {\tiny [1.000,~1.000]} & {\tiny [.994,~.995]} & {\tiny [.725,~.735]} & {\tiny [1.000,~1.000]} & {\tiny [.979,~.980]} \\[0.4ex]
        \bottomrule
    \end{tabularx}
    \begin{minipage}{\linewidth}
        \vspace{0.5ex}
        \footnotesize \textit{Notes:} Each cell reports the empirical coverage probability of the corresponding confidence interval, computed across $250$ Monte Carlo repetitions on data generated from the seven DGPs of \cref{sec:dgp}.The cells reports the average over all support points $\theta_r$ on the prespecified grid $\mathcal{B}$. Brackets contain a 95\% percentile-bootstrap confidence interval for the empirical coverage, obtained by resampling the $250$ Monte Carlo repetitions with replacement (9999 bootstrap samples). FKRB denotes the interval proposed by \citet{FoxKimRyanBajari2011}, see \cref{eq:CIFKRB}. Wald is the sieve Wald interval from \citet{ChenPouzo2015}, see \cref{eq:CIChenPouzo}, and QLR is the (studentized) sieve quasi-likelihood ratio interval, see \cref{eq:CIQLR}. The nominal coverage is $1-\alpha$. A star ($^{*}$) marks cells whose confidence interval for the empirical coverage lies entirely above the nominal level, and a dagger ($^{\dagger}$) marks cells whose interval lies entirely below it.
    \end{minipage}
\end{table}

\begin{table}[htbp]
    \centering
    \caption{Empirical coverage of the original FKRB \citep{FoxKimRyanBajari2011}, and sieve Wald and sieve QLR \citep{ChenPouzo2015} confidence intervals for $\mathbb{E}[\beta_1]$, by data generating process, sample size $N$, grid size $R$ and significance level $\alpha$.}
    \label{tab:coverage_Eb1}
    \footnotesize
    \setlength{\tabcolsep}{1.5pt}
    \begin{tabularx}{\textwidth}{ccc *{12}{>{\centering\arraybackslash}X}}
        \toprule
        & & & \multicolumn{3}{c}{Bimodal} & \multicolumn{3}{c}{Tight normal} & \multicolumn{3}{c}{Wide normal} & \multicolumn{3}{c}{Strictly uniform} \\
        \cmidrule(lr){4-6} \cmidrule(lr){7-9} \cmidrule(lr){10-12} \cmidrule(lr){13-15}
        $\alpha$ & $N$ & $R$ & FKRB & Wald & QLR & FKRB & Wald & QLR & FKRB & Wald & QLR & FKRB & Wald & QLR \\
        \midrule
        0.05 & 1{,}000 & 25 & 0.968 & 1.000 & 0.984 & 0.964 & 1.000 & 0.988 & 0.960 & 1.000 & 0.988 & 0.940 & 1.000 & 0.996 \\
        & & & {\tiny [.938\newline .984]} & {\tiny [.985\newline 1.000]} & {\tiny [.960\newline .994]} & {\tiny [.933\newline .981]} & {\tiny [.985\newline 1.000]} & {\tiny [.965\newline .996]} & {\tiny [.928\newline .978]} & {\tiny [.985\newline 1.000]} & {\tiny [.965\newline .996]} & {\tiny [.903\newline .963]} & {\tiny [.985\newline 1.000]} & {\tiny [.978\newline .999]} \\[0.4ex]
        0.05 & 1{,}000 & 225 & 0.932 & 1.000 & 0.988 & 0.972 & 1.000 & 0.992 & 0.952 & 1.000 & 0.984 & 0.960 & 1.000 & 0.988 \\
        & & & {\tiny [.894\newline .957]} & {\tiny [.985\newline 1.000]} & {\tiny [.965\newline .996]} & {\tiny [.943\newline .986]} & {\tiny [.985\newline 1.000]} & {\tiny [.971\newline .998]} & {\tiny [.918\newline .972]} & {\tiny [.985\newline 1.000]} & {\tiny [.960\newline .994]} & {\tiny [.928\newline .978]} & {\tiny [.985\newline 1.000]} & {\tiny [.965\newline .996]} \\[0.4ex]
        0.05 & 10{,}000 & 25 & 0.944 & 1.000 & 0.976 & 0.964 & 1.000 & 0.988 & 0.952 & 1.000 & 0.988 & 0.964 & 1.000 & 0.968 \\
        & & & {\tiny [.908\newline .966]} & {\tiny [.985\newline 1.000]} & {\tiny [.949\newline .989]} & {\tiny [.933\newline .981]} & {\tiny [.985\newline 1.000]} & {\tiny [.965\newline .996]} & {\tiny [.918\newline .972]} & {\tiny [.985\newline 1.000]} & {\tiny [.965\newline .996]} & {\tiny [.933\newline .981]} & {\tiny [.985\newline 1.000]} & {\tiny [.938\newline .984]} \\[0.4ex]
        0.05 & 10{,}000 & 225 & 0.916 & 1.000 & 0.676 & 0.964 & 1.000 & 0.996 & 0.932 & 1.000 & 0.992 & -- & -- & -- \\
        & & & {\tiny [.875\newline .944]} & {\tiny [.985\newline 1.000]} & {\tiny [.616\newline .731]} & {\tiny [.933\newline .981]} & {\tiny [.985\newline 1.000]} & {\tiny [.978\newline .999]} & {\tiny [.894\newline .957]} & {\tiny [.985\newline 1.000]} & {\tiny [.971\newline .998]} &  &  &  \\[0.4ex]
        \cmidrule(l){1-15}
        0.5 & 1{,}000 & 25 & 0.512 & 0.960 & 0.536 & 0.464 & 1.000 & 0.664 & 0.556 & 0.996 & 0.612 & 0.444 & 0.972 & 0.604 \\
        & & & {\tiny [.450\newline .573]} & {\tiny [.928\newline .978]} & {\tiny [.474\newline .597]} & {\tiny [.403\newline .526]} & {\tiny [.985\newline 1.000]} & {\tiny [.603\newline .720]} & {\tiny [.494\newline .616]} & {\tiny [.978\newline .999]} & {\tiny [.550\newline .670]} & {\tiny [.384\newline .506]} & {\tiny [.943\newline .986]} & {\tiny [.542\newline .663]} \\[0.4ex]
        0.5 & 1{,}000 & 225 & 0.416 & 1.000 & 0.572 & 0.504 & 1.000 & 0.640 & 0.544 & 1.000 & 0.652 & 0.460 & 1.000 & 0.632 \\
        & & & {\tiny [.357\newline .478]} & {\tiny [.985\newline 1.000]} & {\tiny [.510\newline .632]} & {\tiny [.442\newline .565]} & {\tiny [.985\newline 1.000]} & {\tiny [.579\newline .697]} & {\tiny [.482\newline .605]} & {\tiny [.985\newline 1.000]} & {\tiny [.591\newline .708]} & {\tiny [.399\newline .522]} & {\tiny [.985\newline 1.000]} & {\tiny [.571\newline .689]} \\[0.4ex]
        0.5 & 10{,}000 & 25 & 0.576 & 0.856 & 0.640 & 0.560 & 0.992 & 0.840 & 0.540 & 0.908 & 0.684 & 0.488 & 0.836 & 0.544 \\
        & & & {\tiny [.514\newline .636]} & {\tiny [.807\newline .894]} & {\tiny [.579\newline .697]} & {\tiny [.498\newline .620]} & {\tiny [.971\newline .998]} & {\tiny [.789\newline .880]} & {\tiny [.478\newline .601]} & {\tiny [.866\newline .938]} & {\tiny [.624\newline .738]} & {\tiny [.427\newline .550]} & {\tiny [.785\newline .877]} & {\tiny [.482\newline .605]} \\[0.4ex]
        0.5 & 10{,}000 & 225 & 0.500 & 1.000 & 0.588 & 0.492 & 1.000 & 0.744 & 0.544 & 1.000 & 0.856 & -- & -- & -- \\
        & & & {\tiny [.438\newline .562]} & {\tiny [.985\newline 1.000]} & {\tiny [.526\newline .647]} & {\tiny [.431\newline .554]} & {\tiny [.985\newline 1.000]} & {\tiny [.686\newline .794]} & {\tiny [.482\newline .605]} & {\tiny [.985\newline 1.000]} & {\tiny [.807\newline .894]} &  &  &  \\[0.4ex]
        \bottomrule
    \end{tabularx}
    \begin{minipage}{\textwidth}
        \vspace{0.5ex}
        \footnotesize \textit{Notes:} Each cell reports the empirical coverage probability of the corresponding confidence interval for the linear functional $\mathbb{E}[\beta_1] = \sum_r \theta_r \beta_1^{(r)}$, computed across $250$ Monte Carlo repetitions of the FKRB estimator on data generated from the seven DGPs of \Cref{sec:dgp}. Brackets contain a 95\% Wilson score confidence interval for the empirical coverage. FKRB denotes the interval proposed by \citet{FoxKimRyanBajari2011}, see \cref{eq:CIfunctional}. Wald is the sieve Wald interval from \citet{ChenPouzo2015}, see \cref{eq:CIChenPouzo}, and QLR is the (studentized) sieve quasi-likelihood ratio interval, see \cref{eq:CIQLR}. The nominal coverage is $1-\alpha$; coverage above this level reflects conservativeness.
    \end{minipage}
\end{table}

\end{landscape}
